\section{Overview}\label{sec:overview}

\paragraph{What the plugin is for.} A \emph{labelled transition system}
(LTS) is a set of states and labelled steps between them. In Rocq, the
plugin takes an LTS as an inductive predicate
\code{step : M -> option A -> M -> Prop} without parameters: a step
\code{step s l s'} goes from \code{s} to \code{s'} with label \code{l}, and
the label \code{None} is a \emph{silent} step. Given two LTSs and a start
state in each, the plugin answers four kinds of goal
(\file{theories/Bisimilarity.v}): \code{strong_sim} and
\code{strong_bisimilar} (every step is answered by one step with the same
label), and \code{weak_sim} and \code{weak_bisimilar} (silent steps may be
taken freely around the answer). The tactics are \code{sim}, \code{bisim},
\code{weak_sim} and \code{weak_bisim}.

\paragraph{How a goal is decided and certified.} In three stages.
\begin{enumerate}
\item \textbf{Deciding} (\cref{sec:symbolic}). The engine reads the LTS's
  constructors as rules, reads each state as a \emph{container} (a shape,
  and a value at each position: \cref{sec:containers}), and stores sets of
  states as \emph{list decision diagrams} in Sylvan, a C library. It learns
  each rule's steps on the positions the rule reads, explores both systems,
  and plays the simulation game on pairs of states to its greatest
  fixpoint. The verdict is whether the start pair is in it.
\item \textbf{The certificate} (\cref{sec:prooflog}). The engine
  checks the game's result once more, as images, with Sylvan's \emph{proof log} on: every move from
  the relation found lands among the pairs that can answer it. Every
  result Sylvan computes (a union, a difference, an image) is recorded
  with the case its algorithm took. That log, kept in memory (or in a file
  next to the file being checked, on request), is the certificate: what
  Sylvan computed, step by step.
\item \textbf{The proof} (\cref{sec:prooflog:facts}, \cref{sec:qed}). The
  plugin turns each record into a \emph{fact}: a lemma of
  \file{theories/DdTree.v} applied to the record's nodes. From the facts
  and from proofs generated for the goal's own systems (each state as a
  vector, each learned step a real step, every real step learned), it
  builds the proof: its pieces as constants the kernel checks as they are
  made, and one term that names them, with which the tactic closes the
  goal. Rocq's kernel checks that term at \code{Qed}: it typechecks it.
  For a large goal, the command \code{MeBi Prove} builds the same proof
  with every part of it a constant, each checked as it is declared, and
  declares the lemma (\cref{sec:qed:prove}): it needs less memory.
\end{enumerate}

\paragraph{Reading path.} For a newcomer: \cref{sec:prooflog} first (it
defines every term it uses, with worked examples), then \cref{sec:qed}
for how the goal's proof is built from the certificate (it defines what
it needs from Rocq and recalls the terms of \cref{sec:prooflog}), then
\cref{sec:symbolic} for the engine and \cref{sec:containers} for the
encoding of states, then \cref{sec:optimisations} for how the
optimisations fit the certificate, and \cref{sec:commands,sec:tests} to
run things.
